\section{Auction-Based Multi-Policy Framework}

We formalize the policy synthesis problem for a fixed finite horizon, and then introduce our auction-based framework for decentralized policy composition.

\textbf{Multi-objective MDPs.}
We use multi-objective Markov decision processes (\momdps) to model uncertain environments with multiple reward objectives.
An \momdp with \(m\in \N_{>0}\) objectives is a tuple \(\mathcal{M}=(S,A,T,\mu,R,H)\), where \(S\) is the state space, \(A\) is the action space, \(H\in \N_{>0}\) is the horizon, \(T:S\times A\to \dist(S)\) is the transition kernel, \(\mu\in \dist(S)\) is the initial state distribution, and \(R=\{r_i:S\times A\to \R\}_{i\in [1;m]}\) is the set of reward functions.
A Markov policy is a sequence \(\pi=\{\pi_t\}_{t=0}^{H-1}\), where each \(\pi_t:S\to \dist(A)\) maps the current state to a distribution over actions.
For a trace \(\rho=s^0a^0s^1a^1\ldots s^H\), the reward of objective \(i\) is \(w_i(\rho)=\sum_{t=0}^{H-1} r_i(s^t,a^t)\), and the aggregate reward is \(w(\rho)=\sum_{i=1}^m w_i(\rho)\).
The value of a policy is \(V^\pi=\E^\pi[w(\rho)]\), and an optimal centralized policy is any policy \(\pi^*\in \arg\max_\pi V^\pi\).

We now replace the centralized policy by \(m\) local policies, one for each objective.
At every bidding round, each local policy proposes both an action and a bid; the bid determines which proposed action controls the environment for the next \(\tau\in\N_{>0}\) steps.
We use a continuous bid space \([0,\beta]\), where \(\beta\in\R_{>0}\), and resolve bidding ties uniformly at random.
To discourage overbidding, bids are penalized with coefficient \(\rho\in(0,1)\).
In the \poorman setting, only the selected policy pays its bid penalty; in the \allpay setting, every policy pays its own bid penalty.
Without loss of generality, we assume that \(H\) is a multiple of \(\tau\), since the horizon can be padded with zero-reward terminal transitions.
The following Markov game makes this composition rule explicit.

\begin{definition}[Markov bidding game]\label{def:markov-bidding-game}
	Let \(\mathcal{M}=(S,A,T,\mu,R,H)\) be an \momdp with \(m\) objectives, let \(\tau\in\N_{>0}\) be the bidding interval, let \(\beta\in\R_{>0}\) be the bid cap, let \(\rho\in(0,1)\) be the bid penalty coefficient, and let \(\mathsf{P}\in\{\poorman,\allpay\}\) be the penalty model.
	The induced \(m\)-player Markov bidding game is \(\mathcal{G}^{\mathcal{M}}_{\tau,\beta,\rho,\mathsf{P}}=(\hat S,\hat A,\hat T,\hat R,\hat\mu,H)\), where:
	\begin{itemize}
		\item \(\hat S=S\times [1;m]\times [0;\tau-1]\). A state \((s,k,q)\) records the \momdp state \(s\), the current controlling policy \(k\), and the number \(q\) of steps remaining before the next bidding round. When \(q=0\), a new bidding round occurs and the current controller \(k\) is ignored.
		\item \(\hat A=(A\times [0,\beta])^m\). Each policy \(i\) independently chooses an action \(a_i\in A\) and a bid \(b_i\in[0,\beta]\). We write a joint action as \((\vec a,\vec b)=(a_i,b_i)_{i\in[1;m]}\).
		\item Let \(W(\vec b)=\{i\in[1;m]\mid b_i=\max_j b_j\}\). The transition kernel \(\hat T:\hat S\times\hat A\to\dist(\hat S)\) is defined by
		\begin{equation}
			\hat T((s,k,q),(\vec a,\vec b),(s',k',q'))=
			\begin{cases}
				T(s,a_k,s') & q>0 \land k'=k \land q'=q-1,\\
				\frac{1}{\abs{W(\vec b)}}T(s,a_{k'},s') & q=0 \land k'\in W(\vec b)\land q'=\tau-1,\\
				0 & \text{otherwise}.
			\end{cases}
		\end{equation}
		\item For each policy \(i\), the reward is inherited from the \momdp reward \(r_i\), with a bid penalty charged at bidding states:
		\begin{equation}
			\hat r_i((s,k,q),(\vec a,\vec b),(s',k',q'))
			=
			r_i(s,a_{k'})
			-
			\rho b_i\cdot \mathbf{1}\!\left[q=0\land (i=k'\lor \mathsf{P}=\allpay)\right].
		\end{equation}
		\item The initial distribution \(\hat\mu\) satisfies \(\hat\mu(s,k,0)=\mu(s)/m\) for every \(k\in[1;m]\), and \(\hat\mu(s,k,q)=0\) for \(q>0\).
	\end{itemize}
\end{definition}

Thus a local policy is a strategy in the Markov bidding game, which we write as \(\pi_i=(\pi_i^a,\pi_i^b)\), with an action component \(\pi_i^a\) and a bid component \(\pi_i^b\).
Each player optimizes its own cumulative reward \(\sum_{t=0}^{H-1}\hat r_i^t\), while the auction rule determines which local action is executed.
We next analyze the bidding incentives at a single decision state, using continuation values from this game.

\smallskip
\noindent\textbf{Existence of Nash equilibria.}
Continuous bidding makes the Markov bidding game discontinuous at bid ties, since an arbitrarily small bid increase can change the winning policy.
Simon and Zame~\citep{simon1990discontinuous} prove existence for discontinuous games by allowing the sharing rule at discontinuity points to be part of the equilibrium object.
This does not directly imply existence for our game with a fixed uniform tie-breaking rule.
Nevertheless, in the continuous all-pay equilibrium characterized below, bid ties occur with probability zero, so the particular tie-breaking rule is payoff-irrelevant for that equilibrium.
The winner-pays result below should instead be read conditionally on the existence or selection of a pure equilibrium.

\smallskip
\noindent\textbf{Equilibrium continuation values.}
We consider local policies \(\pi_i=(\pi_i^a,\pi_i^b)\) that form a Nash equilibrium of the Markov bidding game.
Because the horizon is finite, the corresponding equilibrium continuation values can be defined by backward induction.
Set \(V^*_{i,H}(s)=0\) for every objective \(i\) and terminal state \(s\).
Suppose the game is at state \(s\) and bidding time \(t\in\{0,\tau,2\tau,\ldots,H-\tau\}\).

For a policy \(i\), define its winning continuation value by
\begin{equation}
	\widetilde V^{\win}_{i,t}(s)
	\coloneqq
	\max_{\pi_i^a}
	\E\left[
		\sum_{q=0}^{\tau-1} r_i(s^{t+q},a^{t+q})
		+V^*_{i,t+\tau}(s^{t+\tau})
	\right],
\end{equation}
where the action sequence is generated by policy \(i\) during the control window.
Similarly, define the losing continuation value by
\begin{equation}
	\widetilde V^{\lose}_{i,t}(s)
	\coloneqq
	\min_{j\neq i}
	\E\left[
		\sum_{q=0}^{\tau-1} r_i(s^{t+q},a^{t+q})
		+V^*_{i,t+\tau}(s^{t+\tau})
	\right],
\end{equation}
where policy \(j\) controls the action sequence during the window using an action policy that maximizes its own winning continuation value.
The control gain of policy \(i\) at \((s,t)\) is
\begin{equation}
	G_{i,t}(s)
	\coloneqq
	\widetilde V^{\win}_{i,t}(s)-\widetilde V^{\lose}_{i,t}(s).
\end{equation}
This quantity measures the value of winning control for the next bidding interval rather than losing it.

Given bids \(b=(b_1,\ldots,b_m)\), let
\begin{equation}
	W(b)\coloneqq \{k\in [1;m]\mid b_k=\max_{\ell\in[1;m]}b_\ell\}
\end{equation}
be the set of highest bidders, and let \(p_i(b)=\mathbf{1}[i\in W(b)]/\abs{W(b)}\) be the probability that \(i\) wins after tie-breaking.
For a penalty model \(\mathsf{P}\in\{\poorman,\allpay\}\), the one-round bidding payoff is
\begin{equation}
	U^{\mathsf{P}}_{i,t}(s,b)
	=
	\begin{cases}
		\widetilde V^{\lose}_{i,t}(s)+p_i(b)\left(G_{i,t}(s)-\rho b_i\right) & \mathsf{P}=\poorman,\\
		\widetilde V^{\lose}_{i,t}(s)+p_i(b)G_{i,t}(s)-\rho b_i & \mathsf{P}=\allpay.
	\end{cases}
\end{equation}
In the \poorman case, the bid cost is paid only in expectation over winning; in the \allpay case, the bid cost is paid regardless of the winner.

A mixed bidding policy for player \(i\) at \((s,t)\) is the bid component \(\pi^b_{i,t}(\cdot\mid s)\in\dist([0,\beta])\), and a mixed bidding profile is \(\pi^b_t(\cdot\mid s)=(\pi^b_{1,t}(\cdot\mid s),\ldots,\pi^b_{m,t}(\cdot\mid s))\).
We write \(b\sim\pi^b_t(\cdot\mid s)\) when bids are sampled independently as \(b_i\sim\pi^b_{i,t}(\cdot\mid s)\).
A mixed bidding profile \(\pi^{b*}_t(\cdot\mid s)\) is a Nash equilibrium of the one-round bidding game at \((s,t)\) if, for every player \(i\) and every alternative bid distribution \(\pi^b_{i,t}(\cdot\mid s)\in\dist([0,\beta])\),
\begin{equation}
	\E_{b\sim\pi^{b*}_t(\cdot\mid s)}\!\left[U^{\mathsf{P}}_{i,t}(s,b)\right]
	\geq
	\E_{b_i\sim\pi^b_{i,t}(\cdot\mid s),\; b_{-i}\sim\pi^{b*}_{-i,t}(\cdot\mid s)}\!\left[U^{\mathsf{P}}_{i,t}(s,b_i,b_{-i})\right].
\end{equation}
The equilibrium continuation value is then
\begin{equation}
	V^*_{i,t}(s)
	=
	\E_{b\sim \pi^{b*}_t(\cdot\mid s)}\left[U^{\mathsf{P}}_{i,t}(s,b)\right],
\end{equation}
which completes the backward-induction definition.

\smallskip
\noindent\textbf{Theoretical analysis of the bidding schemes.}
The first result says that the winner-pays mechanism can implement the largest-gain policy in arbitrarily good pure approximate equilibria.
Recall that a pure bid profile is a pure \(\varepsilon\)-Nash equilibrium if no policy can improve its payoff by more than \(\varepsilon\) through a unilateral bid deviation.

\begin{theorem}[Winner-pays approximate implementation]
	In the \poorman setting, consider a single bidding round at state \(s\) and time \(t\).
	Let \(i^*\) be a policy with largest control gain, let \(j^*\) be a policy with second-largest control gain, and suppose
	\begin{equation}
		G_{i^*,t}(s)>G_{j^*,t}(s).
	\end{equation}
	For any \(\varepsilon>0\) such that \((G_{j^*,t}(s)+\varepsilon)/\rho\leq \beta\) and \(G_{i^*,t}(s)>G_{j^*,t}(s)+\varepsilon\), there is a pure \(\varepsilon\)-Nash equilibrium in which policy \(i^*\) wins control with probability one.
\end{theorem}

\begin{proof}
	Consider the bid profile \(b_{j^*}=G_{j^*,t}(s)/\rho\), \(b_{i^*}=(G_{j^*,t}(s)+\varepsilon)/\rho\), and \(b_k\leq b_{j^*}\) for every \(k\notin\{i^*,j^*\}\).
	Policy \(i^*\) wins uniquely.
	No losing policy can gain by deviating to another losing bid.
	If \(j^*\) deviates to a winning bid, it must bid at least \(b_{i^*}\), giving payoff at most
	\begin{equation}
		\widetilde V^{\lose}_{j^*,t}(s)+G_{j^*,t}(s)-\rho b_{i^*}
		=
		\widetilde V^{\lose}_{j^*,t}(s)-\varepsilon,
	\end{equation}
	which is worse than losing.
	Every other losing policy has control gain at most \(G_{j^*,t}(s)\), so the same argument applies.
	
	Policy \(i^*\) cannot gain by raising its bid.
	Its best deviation is to lower its bid while still winning, which can improve its payoff by at most \(\varepsilon\), since it must bid above \(b_{j^*}\) to remain the unique winner.
	Deviating to a losing bid gives payoff \(\widetilde V^{\lose}_{i^*,t}(s)\), while the proposed bid gives
	\begin{equation}
		\widetilde V^{\lose}_{i^*,t}(s)+G_{i^*,t}(s)-G_{j^*,t}(s)-\varepsilon,
	\end{equation}
	which is strictly larger by assumption.
	Thus no policy can improve by more than \(\varepsilon\), so the profile is a pure \(\varepsilon\)-Nash equilibrium.
\end{proof}

The all-pay mechanism admits a standard mixed-equilibrium construction in the continuous complete-information setting.

\begin{theorem}[All-pay concentration]
	In the \allpay setting, consider a single bidding round at state \(s\) and time \(t\), with continuous bid space \([0,\beta]\).
	Let \(i^*\) and \(j^*\) be policies satisfying
	\begin{equation}
		G_{i^*,t}(s)
		>
		G_{j^*,t}(s)
		>
		\max_{k\notin\{i^*,j^*\}}G_{k,t}(s),
	\end{equation}
	and suppose \(\rho\beta\geq G_{j^*,t}(s)\).
	Then the continuous all-pay bidding game has no pure Nash equilibrium, and there is a mixed Nash equilibrium given by the following bid distributions:
	\begin{itemize}
		\item policy \(i^*\) samples uniformly from \([0,G_{j^*,t}(s)/\rho]\);
		\item policy \(j^*\) places an atom of mass
		\begin{equation}
			\frac{G_{i^*,t}(s)-G_{j^*,t}(s)}
			{G_{i^*,t}(s)}
		\end{equation}
		at \(0\), and otherwise has density \(\rho/G_{i^*,t}(s)\) on \((0,G_{j^*,t}(s)/\rho]\);
		\item every other policy bids \(0\) with probability one.
	\end{itemize}
	In this equilibrium, policy \(i^*\) wins control with probability
	\begin{equation}
		1-\frac{G_{j^*,t}(s)}
		{2G_{i^*,t}(s)}.
	\end{equation}
\end{theorem}

\begin{proof}
	This is a standard complete-information all-pay equilibrium~\citep{baye1996all}, applied to valuations \(G_{i,t}(s)\) and bid cost \(\rho b_i\).
	The bid-cap assumption ensures that the equilibrium support lies inside \([0,\beta]\).
	The stated winning probability follows by integrating the probability that \(j^*\)'s bid is below the bid of \(i^*\):
	\begin{equation}
	\begin{aligned}
		P[i^*\text{ wins}]
		&=
		\int_0^{G_{j^*,t}(s)/\rho}
		\left(
			\frac{G_{i^*,t}(s)-G_{j^*,t}(s)}
			{G_{i^*,t}(s)}
			+
			\frac{\rho b}{G_{i^*,t}(s)}
		\right)
		\frac{\rho}{G_{j^*,t}(s)}
		db\\
		&=
		1-\frac{G_{j^*,t}(s)}
		{2G_{i^*,t}(s)}.
	\end{aligned}
	\end{equation}
\end{proof}

\smallskip
\noindent\textbf{Learning policies in Markov bidding games.}
In our experiments, the transition kernel of the original \momdp is unknown, and therefore the induced Markov bidding game in Definition~\ref{def:markov-bidding-game} is also unknown.
We learn the local policies directly with reinforcement learning.
Following the multi-agent PPO recipe of \citet{yu2022mappo}, we train all local policies concurrently with proximal policy optimization (PPO; \citealp{schulman2017ppo}).
Although the theoretical analysis above uses the continuous bid space \([0,\beta]\), our implementation discretizes this interval into finitely many bid levels and treats the bid as an additional categorical action.
This makes the bidding component compatible with standard policy-gradient implementations, while retaining the interpretation of bids as calibrated control gains.
The continuous-equilibrium statements above therefore do not apply directly to the discrete-bid implementation.
Instead, our experiments evaluate whether PPO converges to useful approximate equilibria under this discretized mechanism, and show that the learned policies reliably coordinate through bids in the environments we study.
Since the objectives in our environments have shared structure, the local policies share the same actor and critic networks.
At each PPO iteration, we collect rollouts from each policy's perspective and sample minibatches from the pooled rollout data to update the shared networks.

To handle a varying number of objectives, we use an attention pooling architecture that compresses the active objective set into a fixed-dimensional embedding.
Each objective information vector is first processed by a fully connected network to produce an objective embedding.
The embeddings are then aggregated through a weighted sum, where the weights are computed from the dot product between each embedding and a learned query vector.
This produces a fixed-dimensional representation of the objective set while allowing the learned policies to adapt to different numbers of active objectives at deployment time.
